为了证明 $\Th{Q}$ 能表示所有可计算函数，我们需要一个精确的可计算性模型作为证明的基础。
当然，存在许多可计算性模型，例如 Turing（图灵）机。
历史上起重要作用的模型之一——它是最早被提出的模型之一，也是 G\"odel 本人所使用的模型——是所谓的\textbf{递归函数}。
递归函数是一类算术函数——也就是说，它们的定义域和值域都是自然数——它们可以由几个基本函数通过几种操作来定义。
基本函数包括常零函数~$\Zero$（$\Zero(x) = 0$）、直接后继函数~$\Succ$（$\Succ(x) = x+1$）以及投影函数。
操作包括\textbf{复合}、\textbf{原始递归}和\textbf{正则极小化}。
复合简单来说就是函数``链接''的一般形式：先应用一个函数，再对所得结果应用另一个函数。
原始递归从两个已定义的函数 $g$、$h$ 出发定义新函数~$f$，规定 $f$ 在 $0$ 处的值由~$g$ 给出，而在任意数 $n+1$ 处的值由 $h$ 应用于 $f(n)$ 给出。
仅用这两个原理定义的函数称为\textbf{原始递归的}。
一个关系是原始递归的，当且仅当其示性函数是原始递归的。
事实上，一大批有趣的函数和关系都是原始递归的（例如加法、乘法、指数运算、整除性），并且我们可以用诸如有界量化和有界极小化等原理，从已有的原始递归函数和关系定义出新的函数和关系。
特别地，这使得我们能够证明，可以用原始递归的方式处理数的\emph{序列}。
也就是说，存在一种将数列``编码''为单个数的方法，使得我们可以利用作用在这些编码上的原始递归函数来计算第 $i$ 个元素、序列的长度、两个序列的连接等等。
为了得到所有可计算函数，我们最终在复合和原始递归的基础上增加了通过\textbf{正则极小化}进行的定义。
函数~$g(x, y)$ 是\textbf{正则}的，当且仅当对每个~$y$，它至少对一个~$x$ 取值为~$0$。
如果 $g$ 是正则的，那么使得 $g(x, y) = 0$ 的最小 $x$ 总是存在的，并且可以通过计算 $g(0, y)$、$g(1, y)$ 等等直到其中一个等于 $0$ 来找到。
由此得到的函数~$f(y) = \umin{x}{g(x, y) = 0}$ 就是从 $g$ 通过正则极小化定义的函数。
它总是全函数且可计算。
由此得到的函数称为\textbf{一般递归的}。
Church（邱奇）-Turing 论题的一个版本指出，可计算算术函数恰好就是一般递归函数。